\section{Discussion}
Results show that PCR-induced chimeras in short-read mitochondrial data can be detected and effectively removed using interpretable features derived from standard alignments, and that a sequence-based CNN1D can further improve sensitivity.

The observed assembly behavior is consistent with the classification profiles of the two models. Gradient Boosting achieved higher precision but lower recall, indicating a more conservative filtering strategy that removed fewer clean reads at the cost of leaving more chimeric reads behind. This explains why Gradient Boosting performed more robustly at low read depth, where preserving true reads is critical for maintaining assembly support. In contrast, CNN1D achieved much higher recall but lower precision, indicating a more aggressive filtering strategy that removed more chimeric reads but also misclassified more clean reads as chimeric. This trade-off appears advantageous at high read depth, particularly in the 20K-read, 50\% chimera condition, where CNN1D fully restored a 1-node graph while Gradient Boosting left residual branching. However, at 3200 reads, the lower precision of CNN1D resulted in over-filtering, removing enough true reads to fragment the graph despite improved chimera detection. Because the present evaluation used simulated reads from a single mitochondrial reference, performance on empirical datasets may differ. Real sequencing libraries can contain platform-specific errors, uneven coverage, nuclear mitochondrial segments, contaminants, and library-preparation artifacts that were not fully represented in the simulation. Different taxa may also vary in mitochondrial genome structure, repeat content, GC composition, and divergence from the selected reference, all of which may affect alignment patterns and sequence-discontinuity features. Likewise, library preparation protocols that use different polymerases, cycle numbers, or insert-size distributions may produce different chimera profiles. These factors may change the precision--recall balance of both models and may require threshold calibration before deployment on new datasets.